\section{Introduction}
The dynamic algorithms paradigm is becoming increasingly popular for studying algorithmic questions in the presence of gradually changing inputs.
A natural goal in this area is to design algorithms that process each change to the input as fast as possible to adapt the algorithm's output (or a data structure for querying the output) to the current state of the input.
The time spent after an update to perform these computations is called the \emph{update time} of the algorithm.
In many cases, bounds on the update time are obtained in an \emph{amortized} sense as an average over a long enough sequence of updates. Among dynamic graph problems, the question of maintaining exact or approximate shortest paths has received considerable attention in the past two decades. The main focus usually lies on maintaining a distance oracle that answers queries for the distance between a pair of nodes.
For this problem, we call an algorithm \emph{fully dynamic} if it supports both insertions and deletions of edges, and \emph{partially dynamic} if it supports only one type of updates; in particular we call it \emph{decremental} if it only supports edge deletions  (which is the focus of this paper), and \emph{incremental} if it only supports edge insertions. 

The running times of partially dynamic algorithms are usually characterized by their bounds on the \emph{total update time}, which is the accumulated time for processing all updates in a sequence of at most $ m $ deletions (where $ m $ is the maximum number of edges ever contained the graph).
A typical design choice, which we also impose in this paper, is small (say polylogarithmic) query time. In particular, our algorithms will have constant query time.
While fully dynamic algorithms are more general, the restriction to only one type of updates in partially dynamic algorithms often admits much faster update times.
In particular, some partially dynamic algorithms have a total update time that almost matches the running time of the fastest static algorithm, i.e., computing \emph{all} updates does not take significantly more time than processing the graph once.

\paragraph{Decremental shortest paths.} For the decremental single-source shortest paths problem, conditional lower bounds~\cite{RodittyZ11,HenzingerKNS15} suggest that \emph{exact} decremental algorithms have an $\Omega(mn)$ bottleneck in their total update time (up to subpolynomial factors).
On the other hand, this problem admits a $ (1 + \epsilon) $-approximation (also called \emph{stretch}) with total update time $ m^{1 + o(1)} $ in weighted, undirected graphs~\cite{BR11,henzinger2014decremental,BernsteinC16,Bernstein17,BernsteinGS21}, which exceeds the running time of the state-of-the art static algorithm by only a subpolynomial factor. Hence for the single-source shortest paths problem on undirected graphs we can fully characterize for which multiplicative stretches the total update time of the fastest decremental algorithm matches (up to subpolynomial factors) the running time of the fastest static algorithm.

Obtaining a similar characterization for the decremental all-pairs shortest paths (APSP) problem is an intriguing open question.
Conditional lower bounds~\cite{DHZ00,HenzingerKNS15} suggest an $\Omega(mn)$ bottleneck in the total update time of decremental APSP algorithms with (a) any finite stretch on \textit{directed} graphs and (b)  with any stretch guarantee with a multiplicative term of $ \alpha \geq 1 $ and an additive term of $ \beta \geq 0 $ with $ 2 \alpha + \beta < 4 $ on \textit{undirected} graphs. This motivates the study of decremental $ (\alpha, \beta) $-approximate APSP algorithms such that $ 2 \alpha + \beta \geq 4 $, i.e., with multiplicative stretch $ \alpha \geq 2 $ or additive stretch $ \beta \geq 2 $.


Apart from two notable exceptions~\cite{HKN2016,AbrahamC13}, all known decremental APSP algorithms fall into one of two categories.
They either (a) maintain exact distances or have a relatively small multiplicative stretch of $ (1 + \epsilon) $ or (b) they have a stretch of at least~$3$.
The space in between is largely unexplored in the decremental setting.
This stands in sharp contrast to the static setting where for undirected graphs approximations guarantees different from multiplicative $ (1 + \epsilon) $ or ``$ 3 $ and above'' have been the focus of a large body of works~\cite{AingworthCIM99,DHZ00,CohenZ01,BaswanaK10,BermanK07,Kavitha12,BaswanaGS09,PatrascuR14,Sommer16,Knudsen17,AkavR20,AkavR21}. 
The aforementioned exceptions in dynamic algorithms~\cite{HKN2016,AbrahamC13}, concern unweighted undirected graphs and maintain $ (1 + \epsilon, 2) $-approximations that simultaneously have a multiplicative error of~$ 1 + \epsilon $ and an additive error of~$ 2 $, which implies a purely multiplicative $ (2 + \epsilon) $-approximation. The algorithms run in $\tilde O(n^{5/2})$\footnote{In the introduction, we make two simplifying assumptions: (a) $ \epsilon $ is a constant and (b) the ratio $ W $ between the maximum and minimum weight is polynomial in $ n $. Unless otherwise noted, the cited algorithms have constant or polylogarithmic query time. Throughout we use $ \tilde O (\cdot) $ notation to omit factors that are polylogarithmic in $ n $.} and $n^{5/2+o(1)}$ total update time respectively. In the next section, we will detail how our algorithms improve over this total update time and generalize to weighted graphs. 

Concerning the fully dynamic setting, Bernstein~\cite{Bernstein09} provided an algorithm for $(2+\epsilon)$-APSP which takes $m^{1+o(1)}$ time per update. Although we are not aware of any lower bounds, it seems to be hard to beat this: no improvements have been made since. The lack of progress in the fully-dynamic setting motivates the study in a partially dynamic setting, where we obtain improvements for the decremental case.


\subsection{Our Results}
In this paper, we provide novel decremental APSP algorithms with approximation guarantees that previously were mostly unexplored in the decremental setting.
Our algorithms are randomized and we assume an oblivious adversary. 
For each pair of nodes, we can not only provide the distance estimate, but we can also report a shortest path~$\pi$ of this length in $\tilde O(|\pi|)$ time using standard techniques, see e.g.~\cite{ForsterGNS23}. 

\paragraph{$\mathbf{(2+\epsilon)}$-APSP for weighted graphs.} Our first contribution is an algorithm for maintaining a $ (2 + \epsilon) $-approximation in weighted, undirected graphs.

\begin{restatable}{theorem}{mainthm}\label{thm:main_mult}
    Given a weighted graph $G$ and a constant $0 < \epsilon < 1$, there is a decremental data structure that maintains a $(2+\epsilon)$-approximation of APSP. The algorithm has constant query time and the total update time is w.h.p.
    \begin{itemize}
        \item $\tilde O(m^{1/2}n^{3/2}\log^2(nW))$ if $m=n^{1+\Omega(1)}$ and $m\leq n^{2-\rho}$ for an arbitrary small constant $\rho$,
        \item $\tilde O(m^{1/2}n^{3/2+\rho}\log^2(nW))$ if $m\geq n^{2-\rho}$ for any constant $\rho$,
        \item $\tilde O(m^{1/2}n^{3/2+o(1)}\log^2(nW))$ otherwise,
    \end{itemize} where $W$ is the ratio between the maximum and minimum weight.
\end{restatable}

% \mtodo{I wonder why we need 3 different parts here, it seems that the 2nd one is new ($m \geq n^{2-\rho}$)?}

The fastest known algorithm with an approximation ratio at least as good as ours is the $ (1 + \epsilon) $-approximate decremental APSP algorithm by Bernstein~\cite{Bernstein16} with total update time $ \tilde O (m n) $. 
With our $ (2 + \epsilon) $-approximation, we improve upon this total update time when $ m =n^{1+\eta} $, for any $\eta>0$.
Furthermore, the fastest known algorithm with a larger approximation ratio than our algorithm is the $ (3 + \epsilon) $-approximate decremental distance oracle by Łącki and Nazari~\cite{LN2020} with total update time $ \tilde O (m \sqrt{n}) $. 
Our result also has to be compared to the \emph{fully dynamic} algorithm of Bernstein~\cite{Bernstein09} for maintaining a $ (2 + \epsilon) $-approximation that takes amortized time $ m^{1 + o(1)} $ per update.
Note that in \emph{unweighted} graphs, $ (2 + \epsilon) $-approximate decremental APSP algorithm can be maintained with total update time $ \tilde O (n^{2.5}) $ (which is implied by the results of Abraham and Chechik~\cite{AbrahamC13} and Henzinger, Krinninger, and Nanongkai~\cite{HKN2016}). Our approach improves upon this bound for most densities, and matches it for $m=n^2$.


\paragraph{$\mathbf{(2+\epsilon, 1)}$-APSP for unweighted and $\mathbf{(2+\epsilon, W_{u,v})}$-APSP for weighted graphs.} 
Our second contribution is a faster algorithm with an additional additive error term of $1$ for unweighted graphs, which in turn can be used to obtain our unweighted $(2+\epsilon)$-APSP result.
The corresponding generalization to weighted graphs can be formulated as follows.

\begin{restatable}{theorem}{thmtwoW}\label{thm:(2_1)approx}
    Given a weighted graph $G$ and a constant $0 < \epsilon < 1$, there is a decremental data structure that maintains a $(2+\epsilon,W_{u,v})$-approximation for APSP, where $W_{u,v}$ is the maximum weight on a shortest path from $u$ to $v$. The algorithm has constant query time and the total update time is w.h.p.~$\tilde O( nm^{3/4}\log^2(nW))$ if $m=n^{1+\Omega(1)}$, and $\tilde O( n^{1+o(1)}m^{3/4}\log^2(nW))$ otherwise, where $W$ is the ratio between the maximum and minimum weight. 
\end{restatable}

\paragraph{$\mathbf{(2+\epsilon)}$-APSP for unweighted graphs.}
We obtain this result by a general reduction from mixed approximations to purely multiplicative approximations, which might be of independent interest beyond the dynamic setting. See Theorem~\ref{thm:reduction} for the statement. We can then combine this with Theorem~\ref{thm:(2_1)approx} to obtain a fast algorithm for non-dense unweighted graphs. 

\begin{restatable}{theorem}{thmtwoUnw}\label{thm:unweighted_2}
    Given an undirected unweighted graph $G$, there is a decremental data structure that maintains a $(2+\epsilon)$-approximation for APSP with constant query time. W.h.p.\ the total update time is bounded by $\tilde O( m^{7/4})$ if $m=n^{1+\Omega(1)}$, and $\tilde O( m^{7/4+o(1)})$ otherwise.
\end{restatable}

Note that for $m=o(n^{8/7})$ this gives the first subquadratic decremental $(2+\epsilon)$-approximate APSP algorithm. 
For $m\leq n^{6/5}$, this approach is faster than our weighted result, \Cref{thm:main_mult}, which was already beating the unweighted state-of-the-art~\cite{AbrahamC13,HKN2016}. 


\paragraph{$\mathbf{(1+\epsilon, 2(k-1))}$-APSP for unweighted graphs.}Our fourth contribution is an algorithm for unweighted, undirected graphs that maintains, for any $ k \geq 2 $, a ($1 + \epsilon$,$2(k-1)$)-approximation, i.e., a distance estimate that has a multiplicative error of~$ 1 + \epsilon $ and an additive error of~$ 2(k-1) $.
\begin{restatable}{theorem}{thmadd}\label{thm:main_add}
Given an undirected unweighted graph $G$, a constant $0 < \epsilon < 1$ and an integer $2 \leq k \leq \log{n}$, there is a decremental data structure that maintains $(1+\epsilon, 2(k-1))$-approximation for APSP with constant query time. The expected total update time is bounded by $\min\{O(n^{2-1/k+o(1)}m^{1/k}), \tilde{O}((n^{2-1/k}m^{1/k}) O(1/\epsilon)^{k/\rho})\}$ where $m=n^{1+\rho}$.
\end{restatable}

Note that for small values of $k$ and if $m=n^{1+\rho}$ for a constant $\rho>0$, we get total update time of $\tilde{O}(n^{2-1/k}m^{1/k})$, and otherwise we have an extra $n^{o(1)}$ factor. In addition, in the special case of $k=\log{n}$, we get a near-quadratic update time of $O(n^{2+o(1)})$.
The state-of-the-art for a purely multiplicative $ (1 + \epsilon) $-approximation is the algorithm of Roditty and Zwick with total update time $ \tilde O (m n) $.\footnote{Note that the algorithms of Roditty and Zwick~\cite{RZ12} for unweighted, undirected graphs precedes the more general algorithm of Bernstein~\cite{Bernstein16} for weighted, directed graphs.}
It was shown independently by Abraham and Chechik~\cite{AbrahamC13} and by Henzinger, Krinninger, and Nanongkai~\cite{HKN2016} how to improve upon this total update time bound at the cost of an additional small additive error term: a $ (1 + \epsilon, 2) $-approximation can be maintained with total update time $ \tilde O (n^{2.5}) $.
This has been generalized by Henzinger, Krinninger, and Nanongkai~\cite{HKN14} to an additive error term of $ 2 (1 + \tfrac{2}{\epsilon})^{k-2} $ and total update time $ \tilde O (n^{2 + 1/k} O(\tfrac{1}{\epsilon})^{k-1}) $.
We improve upon this tradeoff in two ways: (1) our additive term is independent of $ 1/\epsilon $ and \emph{linear} in $k$ and (2) our algorithm profits from graphs being sparse.

\paragraph{Static follow-up work.} Interestingly, our techniques inspired a new algorithm for a static $2$-approximate distance oracle~\cite{DoryFKNVV23}, giving the first such distance oracle algorithm with subquadratic construction time for sparse graphs. Moreover, for $m=n$ they match conditional lower bounds~\cite{PatrascuRT12, aboud2022stronger}.

%their bounds match the space, and hence time, lower bounds of~\cite{PatrascuRT12}, which is conditional on the set intersection conjecture. Their time bound also matches the recent lower bounds conditional on the 3-SUM conjecture~\cite{aboud2022stronger}. 


\subsection{Related Work}

\paragraph{Static algorithms.}
The baseline for the static APSP problem are the exact textbook algorithms with running times of $ O (n^3) $ and $ \tilde O (m n) $, respectively.
There are several works obtaining improvements upon these running times by either shaving subpolynomial factors or by employing fast matrix multiplication, which sometimes comes at the cost of a $ (1 + \epsilon) $-approximation instead of an exact result.
See~\cite{Zwick01} and~\cite{Williams18}, and the references therein, for details on these approaches.
For the regime of stretch $ 3 $ and more in undirected graphs, a multitude of algorithms has been developed with the distance oracle of Thorup and Zwick~\cite{TZ2005} arguably being the most well-known constructions.
In the following, we focus on summarizing the state of affairs for approximate APSP with stretch between $ 1 + \epsilon $ and $ 3 $.

In \emph{weighted}, undirected graphs, Cohen and Zwick~\cite{CohenZ01} obtained a $2$-approximation with running time $ \tilde O (m^{1/2} n^{3/2}) $.
This running time has been improved to $ \tilde O (m \sqrt{n} + n^2) $ by Baswana and Kavitha~\cite{BaswanaK10} and, employing fast matrix multiplication, to $ \tilde O (n^{2.25}) $ by Kavitha~\cite{Kavitha12}.
For $(2+\epsilon)$-APSP, Dory et al.~\cite{DoryFKNVV23} provide two algorithms, for dense graphs: $ O(n^{2.214})$ and sparse graphs: $\tilde O(mn^{2/3})$.
In addition, efficient approximation algorithms for stretches of $ \tfrac{7}{3} $~\cite{CohenZ01} and $ \tfrac{5}{2} $~\cite{Kavitha12} have been obtained.
These results have recently been generalized by Akav and Roditty~\cite{AkavR21} who presented an algorithm with stretch $ 2 + \tfrac{k-2}{k} $ and running time $ \tilde O (m^{2/k} n^{2 - 3/k} + n^2) $ for any $ k \geq 2 $.

In \emph{unweighted}, undirected graphs, a $ (1, 2) $-approximation algorithm with running time $ \tilde O (n^{2.5}) $ has been presented by Aingworth, Chekuri, Indyk and Motwani~\cite{AingworthCIM99}. This has been improved by Dor, Halperin, and Zwick~\cite{DHZ00} that showed a $(1,2)$-approximation with running time $\tilde{O}(\min\{n^{3/2}m^{1/2}, n^{7/3} \})$. They also show a generalized version of the algorithm that gives stretch $ (1, k) $ and running time ${\tilde{O}}(\min\{n^{2-{2}/{(k+2)}}m^{{2}/{(k+2)}}, n^{2+{2}/{(3k-2)}}\})$ for every even $k>2$.
Recently faster $(1,2)$-approximation algorithms based on fast matrix multiplication techniques were developed \cite{DengKRWZ22,durr2023improved}, the fastest of them runs in $O(n^{2.260})$ time \cite{durr2023improved}. This was extended to a $(1+\epsilon,2)$-approximation in $O(n^{2.152})$~\cite{DoryFKNVV23}. In addition, recently Roditty \cite{Roditty23} extended the approach of \cite{DHZ00} to obtain a combinatorial $(2,0)$-approximation for APSP in $\tilde{O}(n^{2.25})$ time in unweighted undirected graphs.  Using fast matrix multiplication, this running time was improved to $O(n^{2.032})$~\cite{DoryFKNVV23,saha2024faster}.


Berman and Kasiviswanathan~\cite{BermanK07} showed how to compute a $ (2, 1) $-approximation in time $ \tilde O (n^2) $.
Subsequent works~\cite{BaswanaGS09,BaswanaK10,PatrascuR14,Sommer16,Knudsen17} have improved the polylogarithmic factors in the running time and the space requirements for such nearly $2$-approximations.
Recently, slightly subquadratic algorithms have been given: an algorithm with stretch $ (2 (1 + \epsilon), 5)$ by Akav and Roditty~\cite{AkavR20}, and an algorithm with stretch $(2,3)$ by Chechik and Zang~\cite{ChechikZ22}.




\paragraph{Decremental algorithms.}

The fastest algorithms for maintaining exact APSP under edge deletions have total update time $ \tilde O (n^3) $~\cite{DemetrescuI06,BaswanaHS07,EvaldFGW21}.
There are several algorithms that are more efficient at the cost of returning only an approximate solution.
In particular, a $ (1+\epsilon) $-approximation can be maintained in total time $ \tilde O (m n) $~\cite{RZ12,Bernstein16,KarczmarzL19}.
If additionally, an additive error of $ 2 $ is tolerable, then a $ (1+\epsilon, 2) $-approximation can be maintained in total time $ \tilde O (n^{2.5}) $ in unweighted, undirected graphs~\cite{HKN2016,AbrahamC13}.
Note that such a $ (1+\epsilon, 2) $-approximation directly implies a $ (2 + \epsilon) $-approximation because the only paths of length $ 1 $ are edges between neighboring nodes.
All of these decremental approximation algorithms are randomized and assume an oblivious adversary.
Deterministic algorithms with stretch $ 1 + \epsilon $ exist for unweighted, undirected graphs with running time $ \tilde O (m n) $~\cite{HKN2016}, for weighted, undirected graphs with running time $ \tilde O (m n^{1+o(1)})$~\cite{BernsteinGS21} and for weighted, directed graphs with running time $ \tilde O (n^3) $~\cite{KarczmarzL20}.

Decremental approximate APSP algorithms of larger stretch, namely at least~$ 3 $, have first been studied by Baswana, Hariharan, and Sen~\cite{BaswanaHS03}.
After a series of improvements~\cite{RZ12, BR11, AbrahamCT14, henzinger2014decremental}, the state-of-the-art algorithms of~\cite{Chechik18,LN2020} maintain $ (2k-1)(1 + \epsilon) $-approximate all-pairs shortest paths for any integer $ k \geq 2 $ and $ 0 < \epsilon \leq 1 $ in total update time $ \tilde O ( (m+n^{1+o(1)})n^{1/k}) $ with query time  $O(\log \log (nW))$ and $O(k)$ respectively.
All of these ``larger stretch'' algorithms are randomized and assume an oblivious adversary.

Recently, deterministic algorithms have been developed by Chuzhoy and Saranurak~\cite{ChuzhoyS21} and by Chuzhoy~\cite{Chuzhoy21}.
One tradeoff in the algorithm of Chuzhoy~\cite{Chuzhoy21} for example provides total update time $ \tilde O (m^{1 + \mu}) $ for any constant $ \mu $ and polylogarithmic stretch.
As observed by M\k{a}dry~\cite{Madry10}, decremental approximate APSP algorithms that are deterministic -- or more generally work against an adaptive adversary -- can lead to fast static approximation algorithms for the maximum multicommodity flow problem via the Garg-Könemann-Fleischer framework~\cite{GargK07,Fleischer00}.
The above upper bounds for decremental APSP have recently been contrasted by conditional lower bounds~\cite{AbboudBKZ22} stating that constant stretch cannot be achieved with subpolynomial update and query time under certain hardness assumptions on 3SUM or (static) APSP.


\paragraph{Fully dynamic algorithms.}
The reference point in fully dynamic APSP with subpolynomial query time is the \emph{exact} algorithm of Demetrescu and Italiano~\cite{DemetrescuI04} with update time $ \tilde O (n^2) $ (with log-factor improvements by Thorup~\cite{Thorup04}).
%The $ n^2 $-barrier can be broken either at the cost of approximation with small query time and/or at the cost of update/query time tradeoffs.
% NOTE We could mention algebraic approaches here and trade-offs of the form mn/t vs t^2
For undirected graphs, several fully dynamic distance oracles have been developed.
In particular, Bernstein~\cite{Bernstein09} developed a distance oracle of stretch $ 2 + \epsilon $ (for any given constant $ 0 < \epsilon \leq 1 $) and update time $ O (m^{1 + o(1)}) $.
In the regime of stretch at least~$ 3 $, tradeoffs between stretch and update time have been developed by Abraham, Chechik, and Talwar~\cite{AbrahamCT14}, and by Forster, Goranci, and Henzinger~\cite{ForsterGH21}.
Finally, most fully dynamic algorithm with update time sensitive to the edge density can be combined with a fully dynamic spanner algorithm leading to faster update time at the cost of a multiplicative increase in the stretch, see~\cite{BaswanaKS12} for the seminal work on fully dynamic spanners.
